\begin{tikzpicture}[
  node distance=1.15cm,
  every node/.style={font=\small},
  flowstep/.style={
    reportprocess,
    minimum width=6.8cm,
    minimum height=0.85cm,
    align=left,
    text width=6.4cm
  },
  flowdecision/.style={
    reportdecision,
    minimum width=1.4cm,
    minimum height=0.7cm,
    align=center,
    inner sep=1pt
  },
  flowerror/.style={
    reportnode,
    minimum width=2.8cm,
    minimum height=0.75cm,
    align=center
  }
]
  \node[flowstep] (s1) {1. 用户在翻译页输入文本并触发分析};
  \node[flowdecision, below=of s1, xshift=-2.2cm] (d1) {输入\\校验};
  \node[flowerror, below=of s1, xshift=2.2cm] (e1) {返回错误};
  \node[flowstep, below=of d1, xshift=2.2cm] (s2) {2. 读取全局关系标识并注入请求体};
  \node[flowstep, below=of s2] (s3) {3. 调用 POST /api/decode};
  \node[flowdecision, below=of s3, xshift=-2.2cm] (d2) {限流\\校验};
  \node[flowerror, below=of s3, xshift=2.2cm] (e2) {返回 429};
  \node[flowstep, below=of d2, xshift=2.2cm] (s4) {4. 执行脱敏、上下文拼接与分析};
  \node[flowdecision, below=of s4, xshift=-2.2cm] (d3) {LLM\\裁决};
  \node[flowstep, below=of s4, xshift=2.2cm] (s5) {5. 启发式降级};
  \node[flowstep, below=of d3, xshift=2.2cm] (s6) {6. 前端渲染结构化结果};

  \draw[reportarrow] (s1) -- (d1);
  \draw[reportarrow] (d1) -- node[midway, left, font=\tiny] {拒绝} (e1);
  \draw[reportarrow] (d1) -- node[midway, above, font=\tiny] {通过} (s2);
  \draw[reportarrow] (s2) -- (s3);
  \draw[reportarrow] (s3) -- (d2);
  \draw[reportarrow] (d2) -- node[midway, left, font=\tiny] {拒绝} (e2);
  \draw[reportarrow] (d2) -- node[midway, above, font=\tiny] {通过} (s4);
  \draw[reportarrow] (s4) -- (d3);
  \draw[reportarrow] (d3) -- node[midway, above, font=\tiny] {降级} (s5);
  \draw[reportarrow] (d3) -- node[midway, left, font=\tiny] {成功} (s6);
  \draw[reportarrow] (s5) -- (s6);
\end{tikzpicture}
